\documentclass[orivec,runningheads]{llncs}
\usepackage[T1]{fontenc}
\usepackage{graphicx}
\usepackage{booktabs}
\usepackage[misc]{ifsym}
\usepackage{xcolor}
\newcommand{\corr}{(\Letter)}
% N.B.: do not change anything above this line. If you require additional packages, please load them directly after this line.

\usepackage{amsmath}
\usepackage[utf8]{inputenc} % allow utf-8 input
\usepackage[colorlinks=true, allcolors=blue]{hyperref}       % hyperlinks
\usepackage{url}            % simple URL typesetting
\usepackage{amsfonts}       % blackboard math symbols
\usepackage{nicefrac}       % compact symbols for 1/2, etc.
\usepackage{microtype}      % microtypography
\usepackage{xcolor}         % colors
\usepackage{graphicx}
\usepackage{float}
\usepackage{bm}
\usepackage{bbm}
\usepackage{enumitem}
\usepackage[colorinlistoftodos]{todonotes}
\usepackage{mwe}
% N.B.: you may delete the preceding line. It is used to display an example image in this template.

% --- REQUIRED PACKAGES FOR TABLE FORMATTING ---
\usepackage{longtable}
\usepackage{booktabs}
\usepackage{array}
\usepackage{hyperref}

\begin{document}

\title{Waste Collection Routing Simulator}

\titlerunning{Waste Collection Routing Simulator}
\authorrunning{A.C. Fernandes et al.}
% First names are abbreviated in the running head.
% If there is one author, write 'A.L. Benjamin'.
% If there are two authors, write 'A.L. Benjamin and C.C. Broadus Jr.'
% If there are more than two authors, '[...] et al.' is used.

\institute{INESC-ID, Lisboa, Portugal \email{afonso.fernandes@inesc-id.pt}
\and
Instituto Superior T\'ecnico, Universidade de Lisboa, Lisboa, Portugal \email{afonso.fernandes@tecnico.ulisboa.pt}
}
%\email{\{a.l.benjamin,a.a.patton\}@fsu.fake}}

\maketitle              % typeset the header of the contribution

\begin{abstract}
Urban waste collection is a critical municipal service with profound implications for sanitation and public health. This service heavily relies on vehicle routing and has traditionally been modeled as a static Capacitated Vehicle Routing Problem (CVRP), where vehicles perform an exhaustive collection of all bins. The introduction of smart waste bin monitoring technology has enabled dynamic vehicle routing, allowing vehicles to selectively visit only the subset of bins that require immediate service or are profitable to collect, a task which can be modeled as a Vehicle Routing Problem with Profits (VRPP). Furthermore, this task must be performed across an extended planning horizon, where routes executed in a given period impact the optimal routes for subsequent periods, forming a Multi-Period VRPP (MPVRPP). Because of the large scale and temporal dependencies, transferring existing CVRP solvers to this selective routing problem presents significant structural difficulties. Moreover, progress is currently hindered by a lack of standardized multi-period evaluation environments. To address these limitations, this paper introduces WSmartRoute+, a Python simulation framework designed specifically for the stochastic MPVRPP. Using this environment, a unified methodology is developed to adapt state-of-the-art CVRP algorithms - including Hybrid Genetic Search (HGS), Adaptive Large Neighborhood Search (ALNS), and recent Neural Combinatorial Optimization (NCO) models - to dynamic waste collection. The adaptations introduce novel profit-aware repair operators, and dynamic dispatch heuristics. To evaluate the generality of the proposed adaptations, an extensive benchmark of these classical and learning-based solvers is conducted. The experimental results establish a unified baseline for the MPVRPP and demonstrate WSmartRoute+ as a testbed for future research in dynamic, multi-period waste collection routing.
\end{abstract}

\begin{center}
\scriptsize
\renewcommand{\arraystretch}{1.3}
\begin{longtable}{
  >{\raggedright\arraybackslash}p{0.14\textwidth}
  >{\raggedright\arraybackslash}p{0.14\textwidth}
  >{\raggedright\arraybackslash}p{0.12\textwidth}
  >{\raggedright\arraybackslash}p{0.18\textwidth}
  >{\raggedright\arraybackslash}p{0.16\textwidth}
  >{\raggedright\arraybackslash}p{0.16\textwidth}
}
\caption{Taxonomic Deconstruction of Optimization Algorithms: Classification and Search Logic. Detailed descriptions and critiques are provided in Appendix A.} \label{tab:main_taxonomy} \\
\toprule
\textbf{Algorithm} & \textbf{Origin Ref.} & \textbf{Basis Archetype} & \textbf{Classification \& Paradigm} & \textbf{State \& Target Domain} & \textbf{Operators / Search Logic} \\
\midrule
\endfirsthead

\multicolumn{6}{c}%
{{\bfseries \tablename\ \thetable{} -- continued from previous page}} \\
\toprule
\textbf{Algorithm} & \textbf{Origin Ref.} & \textbf{Basis Archetype} & \textbf{Classification \& Paradigm} & \textbf{State \& Target Domain} & \textbf{Operators / Search Logic} \\
\midrule
\endhead

\midrule \multicolumn{6}{r}{{Continued on next page}} \\
\endfoot

\bottomrule
\endlastfoot

\section{Introduction}

\textcolor{blue}{Urban waste collection is a core public service with significant economic, environmental, and social implications. The operational planning of waste collection routes directly affects fuel consumption, labor costs, greenhouse gas emissions, service reliability, and citizen satisfaction. From a methodological standpoint, waste collection routing represents a particularly challenging class of vehicle routing problems, characterized by spatially distributed service points, capacity constraints, heterogeneous waste generation rates, and strict operational requirements.
}

\textcolor{blue}{Over the past decades, Operations Research (OR) has produced a wide range of optimization-based models and solution methods for waste collection routing, including exact formulations, decomposition approaches, and advanced heuristics \cite{hess2024waste,jorge2022hybrid,lopes2023efficient,de2024data}. While these approaches are capable of modeling realistic constraints and multiple objectives, their evaluation is typically conducted either on limited real instances or on static snapshots of the system. As a result, important operational aspects—such as daily variability in waste accumulation, interdependence between routing decisions over time, and the long-term effects of routing policies—are often neglected.}

\textcolor{blue}{In parallel, recent advances in Deep Learning (DL) have given rise to Neural Combinatorial Optimization (NCO) methods, which aim to learn routing policies directly from data using deep neural networks trained via reinforcement learning \cite{wu2024neural}. These approaches have shown promising results on classical routing benchmarks, particularly when evaluated on synthetic and static problem instances. However, their applicability to real-world waste collection operations remains largely unexplored, due in part to the lack of realistic, dynamic, and reproducible evaluation environments.}


\textcolor{blue}{A critical limitation of the current literature is the absence of simulation frameworks that faithfully capture the operational dynamics of waste collection systems. Real-world operations involve continuous waste accumulation, sensor uncertainty, service-level requirements, and feedback effects whereby routing decisions made on one day influence the state of the system on subsequent days. Without a simulation layer that explicitly models these dynamics, it is difficult to assess the true performance, robustness, and scalability of both OR-based solvers and learning-based routing policies.}


\textcolor{blue}{To address this gap, this paper introduces a Waste Collection Routing Simulator designed to serve as a unified experimental platform for the development, evaluation, and comparison of routing strategies. The simulator integrates a dynamic system state that evolves over time, daily route execution, container-level waste accumulation processes, and multi-objective performance metrics reflecting real operational trade-offs. By decoupling the simulation of the operational environment from the optimization or learning components, the framework enables a fair and reproducible comparison between expert-crafted OR solvers and data-driven NCO models under identical conditions.}

\textcolor{blue}{Within this simulator, routing decisions are not evaluated in isolation but as part of a sequential decision-making process, where each day's routes influence future system states, service feasibility, and performance outcomes. This perspective shifts the focus from solving single-shot routing problems to evaluating routing policies over extended planning horizons, which is more aligned with real-world waste management practice.}

\textcolor{blue}{The proposed simulator is validated using real-world data from a waste management company operating across multiple municipalities in Portugal, involving instances with several hundreds of containers. It provides a realistic testbed for studying the strengths and limitations of both optimization-based and learning-based routing approaches, and supports systematic experimentation on scalability, robustness to uncertainty, and multi-objective trade-offs in waste collection operations.}

\section{Related Work}
\subsection{Neural Combinatorial Optimization}
Following a slightly modified version of the taxonomy introduced by \cite{wu2024neural}, which also considers the categories defined by \cite{sun2023difusco}, NCO solvers can be classified into five different types based on how the DNN model generates the solution: 1) autoregressive constructive solvers \cite{Kool2018AttentionLT,kwon2020pomo,vinyals2015pointer}, which start from an empty set and build a valid solution by sequentially selecting unvisited vertices; 2) non-autoregressive constructive solvers \cite{joshi2019efficient,karalias2020erdos}, where the DNN model learns to output the entire solution (usually consisting of probabilities over the vertices or edges of the graph, that are then used by a decoding algorithm to build a valid route) in a single step; 3) improvement heuristic solvers \cite{ma2024learning}, which use a Markov Decision Process (MDP) to iteratively refine an existing complete solution; 4) predict-once solvers \cite{sun2023difusco}, where the model predicts information before starting a solution search with a OR solver or a search algorithm (like beam search); and 5) predict-many solvers \cite{liu2024evolution}, a similar approach to the previous category, with the difference that the model performs additional predictions to further restrict the solution search space, and thus decrease time complexity. The NCO solvers from different categories provide various trade-offs between performance, solving time, and ability to scale to larger problem instances \cite{joshi2020learning,sun2023difusco,wu2024neural}.\\
While all of these algorithms share the same purpose, to produce a near-optimal solution to a given CO problem instance, in practice they learn to perform different specific tasks using a plethora of methods from multiple ML paradigms, including training to mimic OR solvers via SL \cite{joshi2019efficient,sun2023difusco,vinyals2015pointer}, approximating the optimal policy by minimizing a cost function via policy gradient RL \cite{Kool2018AttentionLT,kwon2020pomo,ma2024learning}, and using Unsupervised Learning (UL) to make a model define a probability distribution over a graph's vertices that is used at inference time to build a feasible route through the method of conditional expectation \cite{karalias2020erdos}.

\subsection{Multi-Period Vehicle Routing Problem with Profits}
The Vehicle Routing Problem with Profits (VRPP) can be seen as a variation of the Vehicle Routing Problem (VRP), which is in itself an generalization of the TSP where each vertex has a demand and routes must be constructed such that they both start and end at the depot \cite{Kool2018AttentionLT,wu2024neural}. Unlike the base VRP, the VRPP's objective function is the profit, which considers the demand collected in addition to the route travel cost, thus visiting only the subset of nodes that is profitable, instead of always visiting every available node \cite{doi:10.1137/1.9781611973594.ch10}. Additionally, for the Multi-Period variant of the VRPP (MPVRPP), the vehicles must perform multiple collection routes during a given time range, where the current day's optimal route is dependent upon the choices taken at previous days \cite{DBLP:journals/eor/ZhangCCLQ13}.

\section{Problem Definition}
We model the Multi-Period Vehicle Routing Problem with Profits (MPVRPP) as a profit maximizing, multi-period constrained combinatorial optimization problem over a horizon of $D$ days. 

\subsection{Problem State and Dynamics}
Let $\mathcal{G} = (\mathcal{V}, \mathcal{E})$ be a fully connected graph where $\mathcal{V} = \{v_0, v_1, \dots, v_V \}$, with $v_0$ representing the depot and $\mathcal{V} \setminus \{v_0\}$ denoting the set of waste bins located at spatial coordinates $\vec{l}_i \in \mathbb{R}^2$. The state of the environment at day $d \in \{1, \dots, D\}$ is defined by the fill levels of the bins, denoted by the vector $\vec{w}_d = [w_{1, d}, \dots, w_{V, d}]^\top \in [0, \tau]^V$, where $\tau$ is the maximum bin capacity.

The waste $w_{i,d}$ accumulates stochastically according to a daily waste fill value sampled from a bin-specific distribution $\delta_{i,d} \sim P_i$. If a bin $v_i$ is visited on day $d$, its waste is collected and the level is reset; otherwise, waste continues to accumulate up to the bins' maximum capacity $\tau$:
\begin{equation}
    w_{i, d+1} = 
    \begin{cases} 
        \delta_{i, d+1} & \text{if } v_i \in \mathcal{A}_d \\
        \min(w_{i, d} + \delta_{i, d+1}, \tau) & \text{if } v_i \notin \mathcal{A}_d
    \end{cases}
\end{equation}

\subsection{Optimization Objective}
At each day $d$, the agent selects a route $\mathcal{A}_d = (a_0, a_1, \dots, a_{T_d}, a_{T_d+1})$ starting and ending at the depot $v_0$. The constrained optimization problem for a single day can be formulated as finding a route that:
\begin{align}
    \textrm{maximizes} \quad & \mathcal{P} = r_{w} \cdot \sum_{v_i \in \mathcal{A}_d} w_{i, d} - c_{km} \cdot \sum_{t=0}^{T_d} \|\vec{l}_{a_t} - \vec{l}_{a_{t+1}}\|_2 \label{eq:profit_function}\\
    \textrm{subject to} \quad & a_{0} = a_{T_d} \nonumber\\
    & a_t \neq a_{t'} \quad \forall t \neq t' \in \{1, \dots, T_{d-1}\} \nonumber\\
    & \sum_{v_i \in \mathcal{A}_d} w_{i, d} \leq Q \nonumber
\end{align}
where $r_{w}$ is a scalar value for the revenue obtained from each waste kilogram collected, $c_{km}$ is a scalar value corresponding to the cost it takes for a vehicle to travel one kilometer, and $Q$ is the vehicle maximum waste capacity constraint.

Since decisions made on day $d$ affect the state on day $d+1$, we must optimize over the full horizon. The global objective is to find a policy $\bm{\pi}$ that generates a sequence of routes $(\mathcal{A}_1, \dots, \mathcal{A}_D)$ to maximize the expected cumulative profit:
\begin{equation}
    \max_{\bm{\pi}} \mathbb{E} \left[ \sum_{d=1}^D \mathcal{P}(\mathcal{A}_d, \bm{w}_d) \right],
\end{equation}
where $\mathcal{P}$ is the daily route profit function defined in \eqref{eq:profit_function}.

\section{Methodology}

In this work, we solve the Vehicle Routing Problem with Profits (VRPP) using three state-of-the-art metaheuristics: Adaptive Large Neighborhood Search (ALNS), Hybrid Genetic Search (HGS), and Hyper-Heuristic Ant Colony Optimization (HH-ACO). These algorithms were selected for their proven capability to balance exploration and exploitation in complex combinatorial optimization landscapes.

\subsection{Adaptive Large Neighborhood Search (ALNS)}
The Adaptive Large Neighborhood Search (ALNS) extends the Large Neighborhood Search framework by dynamically selecting destroy and repair operators based on their past performance. At each iteration, a \textit{destroy} operator removes $q$ customers from the current solution, and a \textit{repair} operator reinserts them to form a new candidate solution.

For the VRPP, specific operators are designed to manage the selection of profitable nodes. Destroy operators may target nodes with low profit-to-cost ratios, while repair operators prioritize the insertion of high-profit nodes that incur minimal detour costs.

The selection of operator $i$ is governed by a roulette-wheel mechanism, where the probability $P(i)$ depends on its accumulated weight $w_i$. The weights are updated adaptively after each segment of $\varphi$ iterations:
\begin{equation}
    w_{i,t+1} = \lambda w_{i,t} + (1-\lambda) \frac{\pi_i}{\theta_i}
\end{equation}
where $\lambda \in [0, 1]$ is a decay factor, $\pi_i$ is the accumulated score of operator $i$ during the segment, and $\theta_i$ is the number of times it was used. Scores are awarded based on the quality of the new solution (e.g., finding a new global best).

\subsection{Hybrid Genetic Search (HGS)}
The Hybrid Genetic Search (HGS) combines the evolutionary principles of a genetic algorithm with the fine-tuning capabilities of local search. The algorithm maintains a population of solutions, including feasible and infeasible individuals, to navigate the search space effectively. The core components include:
\begin{itemize}
    \item \textbf{Crossover:} An Ordered Crossover (OX) or similar operator generates offspring by inheriting routes from two parent solutions. In the context of VRPP, the crossover must ensure that the offspring respects the profit collection objectives.
    \item \textbf{Local Search (Education):} Each offspring undergoes an intensive local search procedure (e.g., 2-opt, Relocate, Swap) to reach a local optimum.
    \item \textbf{Population Management:} Diversity is strictly maintained by penalizing solutions that are too similar to existing ones, preventing premature convergence.
\end{itemize}
The fitness $F(S)$ of a solution $S$ is a weighted sum of its cost $C(S)$ and a diversity contribution $D(S)$:
\begin{equation}
    F(S) = C(S) + (1 - \frac{|P_{feas}|}{|P|}) \cdot D(S)
\end{equation}
where $|P_{feas}|$ is the number of feasible individuals in the population $P$.

\subsection{Hyper-Heuristic Ant Colony Optimization (HH-ACO)}
While traditional Ant Colony Optimization (ACO) constructs solutions probabilistically based on pheromone trails, the Hyper-Heuristic ACO (HH-ACO) operates at a higher level of abstraction. Instead of directly building routes, the ants construct a sequence of low-level heuristics (LLH) to apply to the problem.

The problem state can be defined by the current solution quality or search stage. The ants probabilistically select heuristics based on a pheromone matrix $\tau_{sh}$ associated with state $s$ and heuristic $h$. The probability of choosing heuristic $h$ in state $s$ is given by:
\begin{equation}
    p_{sh} = \frac{[\tau_{sh}]^\alpha [\eta_{sh}]^\beta}{\sum_{k \in \mathcal{H}} [\tau_{sk}]^\alpha [\eta_{sk}]^\beta}
\end{equation}
where $\eta_{sh}$ is heuristic information indicating the potential efficacy of $h$, and $\alpha, \beta$ control the relative influence of pheromone and heuristic information. This allows the algorithm to learn essentially "which heuristic to apply when," adapting the search strategy specifically to the instance of the VRPP being solved.

\section{Experimental Evaluation}
We evaluate the models on a scenario where, given a statistical distribution, and the coordinates of a depot and multiple recycling bins, a simulator generates an amount of waste to fill each bin for ${d=\{31, 93, 365\}}$ days. Each day, the models are tasked with selecting which bins the recycling vehicle(s) are to collect waste from, a task which can be interpreted as a type of CO problem.

\subsection{Data}
During the evaluation process, we use the following statistical distributions and data sources (alternatives include \cite{brouwer2023comparison}):
\begin{itemize}
    \item Gamma distribution, and files with the coordinates and the waste accumulation rate of ${b_1=\{20, 50, 225\}}$ bins;
    \item Empirical distribution (WSmart+ Bin Analysis module), and two files with the coordinates of $b_2={317}$ bins and how much waste was generated each day for each bin (from 2021-01-12 to 2023-10-23), respectively.
\end{itemize}
We evaluate the models with ${V_{\Gamma}=20, 50, 100, 150, 225}$ bins on the Gamma distribution and ${V_{E}=20, 50, 100, 150, 225, 317}$ bins on the Empirical distribution. When $V_{\Gamma} \in V_1 \lor V_{E}=V_2$, we run a single simulation with bins available in the coordinates data file. Otherwise, we randomly sample $V_{\Gamma} = \{100, 150\}$ or $V_{E} = \{20, 50, 100, 150, 225\}$ bins for $N_s=100$ or $N_s=10$ samples, respectively.\\
Parameters for the Gamma distribution:
\begin{itemize}
    \item Gamma 1: alpha = [5, 5, 5, 5, 5, 10, 10, 10, 10, 10], beta =[5, 2]
        \subitem mean = [25, 10, 50, 20]
        \subitem variance = [125, 20, 250, 40]
    \item Gamma 2: alpha = [2, 2, 2, 2, 2, 6, 6, 6, 6, 6], beta = [6, 4]
        \subitem mean = [12, 8, 36, 24]
        \subitem variance = [72, 32, 216, 96]
    \item Gamma 3: alpha = [1, 1, 1, 1, 1, 3, 3, 3, 3, 3], beta = [8, 6]
        \subitem mean = [8, 6, 24, 18]
        \subitem variance = [64, 36, 192, 108]
\end{itemize}
The daily generated waste data from the WSmart+ Bin Analysis module is pre-processed to compute the waste accumulation rate of each bin.

\subsubsection{Map}
The map of the $V=317$ bins' coordinates can be seen in Figure \ref{fig:map}.
\begin{figure}
    \centering
    \includegraphics[width=0.75\linewidth]{Images/map317.png}
    \caption{Map of the depot and the 317 bins used in experiments.}
    \label{fig:map}
\end{figure}

\subsubsection{Baselines}
We will compare the proposed models with some state-of-the-art algorithms proposed recently, including: \textsc{look\_ahead\_policy} \cite{jorge2022hybrid}; \textsc{policy\_lastminute} \cite{lopes2023efficient}; \textsc{gurobi} \cite{de2024data}; and \textsc{policy\_regular}, where each location is visited regularly with an interval of 3 days, following a route computed by a TSP solver.

\subsection{Results}
\subsubsection{Gamma Distribution (31 days)}
Comparison of the algorithms and models' efficiency (kg/km) and the number of times a bin overflows with waste (overflows) for ${b_{\Gamma}=20, 50, 100, 150, 225}$ \ref{fig:log_gamma31}.
\begin{figure}
    \centering
    \includegraphics[width=0.7\textwidth]{Images/Results/Gamma/31days/log_comp.png}
    \caption{Experimental results for Gamma distribution (with logarithmic scale on the Y axis and normalized by the number of days).}
    \label{fig:log_gamma31}
\end{figure}\\
Simulation runtime (time): \ref{fig:g31_logtime}\\
Number of bin waste collections (ncol): \ref{fig:g31_ncols}
\begin{figure}
    \centering
    \includegraphics[width=0.49\linewidth]{Images/Results/Gamma/31days/log_time.png}
    \includegraphics[width=0.49\linewidth]{Images/Results/Gamma/31days/ncol.png}
    \caption{(Left) Time (in seconds) it took to complete the simulation (with logarithmic scale on the Y axis and normalized by the number of days). (Right) Total number of bins collected during the simulation (normalized by the number of days).}
    \label{fig:g31_logtime}
    \label{fig:g31_ncols}
\end{figure}\\
In terms of truck capacity efficiency (kg/km) and the amount of waste overflows, the \textsc{attention model} seems to be the best option when it comes to smaller graphs, but \textsc{gurobi} eventually overtakes it as the best algorithms as we increase the number of bins in the simulation.\\
The \textsc{attention model} can output a solution to a problem instance in just a few seconds of inference time, making it the fastest algorithm by far. This is to be expected, as the model performs multiple parallel computations on the GPU, while the other algorithms are not parallelizable and must thus run on a single CPU thread. While significantly slower in comparison, both the \textsc{policy\_regular} and the \textsc{policy\_lastminute} can output a solution in a respectable time frame, taking around 1 to 3 minutes for graphs with $b \geq 100$. The \textsc{look\_ahead\_policy} and \textsc{gurobi} are - by a wide margin - the slowest of all algorithms, taking up to 3500x and 17000x more time than the attention model, respectively, to finish the simulation for 225 bins.\\
Overall, the attention model seems to offer the best runtime and performance in small graphs, while \textsc{gurobi} always offers a solid performance and is the superior algorithm for larger graphs with the trade-off of having a significantly higher runtime and computational cost. Also, \textsc{look\_ahead\_policy} seems to offer no significant performance benefits when compared to other algorithms and is the second most costly algorithm, making it a worse choice regardless of the metric.

\subsubsection{Gamma Distribution (365 days)}
Some results for 20 bins: \ref{fig:enter-label}\\
Some tours for 20 bins: \ref{fig:tours}
\begin{figure}
    \centering
    \includegraphics[width=\linewidth]{Images/comp.png}
    \caption{For gamma dist 1 on 365 days.}
    \label{fig:enter-label}
\end{figure}
\begin{figure}
    \centering
    \includegraphics[width=0.49\textwidth]{Images/am_tour.png}
    \includegraphics[width=0.49\textwidth]{Images/lookahead_tour.png}
    \caption{Example tours for Attention Model (left) and LookAhead Policy (right).}
    \label{fig:tours}
\end{figure}

\subsubsection{WSmart+ Bin Analysis (93 days)}
Comparison of the algorithms and models' efficiency (kg/km) and the number of times a bin overflows with waste (overflows) for ${b_{\Gamma}=20, 50, 100, 150, 225, 317}$ \ref{fig:log_emp93}.
\begin{figure}
    \centering
    \includegraphics[width=0.49\textwidth]{Images/Results/Empirical/93days/log_comp.png}
    \includegraphics[width=0.49\textwidth]{Images/Results/Empirical/93days/comp.png}
    \caption{Experimental results for Gamma distribution (with logarithmic scale on both axes and normalized by the number of days).}
    \label{fig:log_emp93}
\end{figure}\\
Simulation runtime (time): \ref{fig:e93_logtime}\\
Number of bin waste collections (ncol): \ref{fig:e93_ncols}
\begin{figure}
    \centering
    \includegraphics[width=0.49\linewidth]{Images/Results/Empirical/93days/log_time.png}
    \includegraphics[width=0.49\linewidth]{Images/Results/Empirical/93days/ncol.png}
    \caption{(Left) Time (in seconds) it took to complete the simulation (with logarithmic scale on both axes and normalized by the number of days). (Right) Total number of bins collected during the simulation (normalized by the number of days).}
    \label{fig:e93_logtime}
    \label{fig:e93_ncols}
\end{figure}\\
Interestingly, the results seem to differ significantly from previous experiments when it comes to truck capacity efficiency (kg/km) and the amount of waste overflows. Compared to the results for the Gamma distribution, \textsc{gurobi} is no longer the most conservative algorithm with regards to the number of overflows and has an higher truck capacity efficiency. Also, the \textsc{policy\_lastminute} algorithm is extremely competitive in this scenario, achieving the best results in medium-sized graphs ($50 \leq b < 225)$.\\
The runtime results follow the same pattern as in the previous experiments, with the \textsc{attention model} being the fastest algorithm, followed by the \textsc{policy\_regular} and the \textsc{policy\_lastminute}, with the slowest of all being \textsc{gurobi}.

\section{Conclusion and Future Work}
Should write some text here after finishing up the Experimental Evaluation.

\begin{credits}
\subsubsection{\ackname} This study was funded by the WSmart Route+ project (grant number BI|2024/549), and carried out at INESC-ID.

\subsubsection{\discintname}
The authors have no competing interests to declare that are
relevant to the content of this article.
\end{credits}
%
% ---- Bibliography ----
%
% BibTeX users should specify bibliography style 'splncs04'.
% References will then be sorted and formatted in the correct style.
%
\bibliographystyle{splncs04}
\bibliography{mybibliography}

\section{Appendix A: Detailed Taxonomic Critique Table}
\label{appendix:taxonomy}

\begin{center}
\small
\renewcommand{\arraystretch}{1.5}
\begin{longtable}{
  >{\raggedright\arraybackslash}p{0.22\textwidth}
  >{\raggedright\arraybackslash}p{0.34\textwidth}
  >{\raggedright\arraybackslash}p{0.34\textwidth}
}
\caption{Comprehensive Structural Descriptions and Critiques of Optimization Paradigms.} \label{tab:appendix_taxonomy} \\
\toprule
\textbf{Algorithm} & \textbf{Structural Description} & \textbf{True Identity / Critique} \\
\midrule
\endfirsthead

\multicolumn{3}{c}%
{{\bfseries \tablename\ \thetable{} -- continued from previous page}} \\
\toprule
\textbf{Algorithm} & \textbf{Structural Description} & \textbf{True Identity / Critique} \\
\midrule
\endhead

\midrule \multicolumn{3}{r}{{Continued on next page}} \\
\endfoot

\bottomrule
\endlastfoot

Adaptive Large Neighborhood Search (ALNS) & Destroys localized portions of a solution and rebuilds using multiple heuristics. & State-of-the-art routing framework; dynamically controls and weights multiple lower-level heuristics. \\ \midrule

Adaptive Kernel Search (AKS) & Dynamically assesses instance difficulty (EASY, NORMAL, HARD) and solves a sequence of restricted MIPs with adaptive kernel sizing. & Powerful general-purpose MIP heuristic; significantly outperforms standard heuristics by offering a flexible quality-time trade-off. \\ \midrule

Ant Colony Optimization (ACO) & Constructive paths based on local heuristic information and global pheromone stigmergy. & Foundational swarm paradigm; rigorous value in probabilistic construction, though sensitive to decay parameters. \\ \midrule

Artificial Bee Colony (ABC) & Three agent types (Employed, Onlooker, Scout) model foraging behavior. & \textbf{Deceptive.} Isomorphic to Differential Evolution mutation and proportional selection; often requires GA injections for discrete parity. \\ \midrule

Attention Model (AM) & Sequentially constructs a valid route by selecting unvisited vertices using a Transformer. & Foundational RL NCO baseline; successfully replaces hard-coded heuristics with learned policies. \\ \midrule

Attention Model with POMO & Forces the AM to start from multiple roots, generating diverse trajectories to exploit graph symmetry. & Highly effective RL training technique; mathematically sound exploitation of problem invariance. \\ \midrule

Branch-and-Bound (BB) & Systematic enumeration of candidate solutions via state-space search, using bounds to "prune" non-optimal sub-trees. & The foundational mathematical procedure for exact optimization; guaranteed global optimality but computationally prohibitive at scale. \\ \midrule

Branch-and-Price-and-Cut (BPC) & Solves LP relaxations via pricing, tightens bounds with cuts, and branches for integrality. & The gold standard for exact multicommodity flow; rigorously proves optimality via simultaneous column and row generation. \\ \midrule

Circuit-SAT GNN & Iteratively predicts variable assignments to restrict the search tree via continuous circuit loss. & Purely mathematical Unsupervised Learning application for iterative state-space reduction. \\ \midrule

Cluster-First Route-Second & Assigns seed customers via an exact Generalized Assignment Problem (GAP), then sequences routes heuristically. & Foundational decomposition paradigm; leverages exact solvers for spatial clustering where heuristics traditionally struggle. \\ \midrule

Deep Policy Dynamic Programming (DPDP) & SL model predicts edge probabilities (heatmaps) once to restrict the state space for exact DP. & Rigorously synthesizes DL probability predictions with exact, mathematically sound DP bounds. \\ \midrule

DeepACO & Uses a GNN trained via RL to output a single non-autoregressive matrix to initialize pheromones. & Excellent modern hybrid; uses DL to intelligently bias canonical swarm mathematics. \\ \midrule

Differential Evolution (DE) & Adds scaled differences between random members to a target base vector. & Highly effective continuous optimizer; the mathematical engine behind many relabeled "swarm" metaphors. \\ \midrule

Efficient Active Search (EAS) & Predicts a baseline policy once, then uses RL at test-time to update weights for the specific instance. & Novel Predict-Once RL application; effectively bridges the gap between inference and local trajectory search. \\ \midrule

Ensemble Move Acceptance & A group decision-making approach where multiple criteria are combined for joint decisions. & Mitigates individual weaknesses of standalone criteria; majority rule (G-VOT) is often the most robust. \\ \midrule

Erdős Goes Neural (EGN) & Unsupervised GNN outputs soft probabilities processed by sequential decoding for guarantees. & Rigorous bridging of DL probabilities with deterministic, discrete mathematical guarantees. \\ \midrule

Evolution Strategy Mu Plus Lambda (ES-MPL) & $\mu$ parents generate $\lambda$ offspring; deterministic elitist selection from the combined pool. & Foundational archetype establishing elitist survival principles for numerical optimization. \\ \midrule

Evolution Strategy Mu, Kappa, Lambda (ES-MKL) & Parents generate offspring; individuals survive for a maximum of $\kappa$ generations before discard. & Foundational extension of ES introducing chronological memory/lifespan bounds to populations. \\ \midrule

Evolution Strategy Mu, Lambda (ES-MCL) & Selection occurs strictly from $\lambda$ offspring, discarding the parent generation. & Foundational archetype introducing the mechanism of escaping local optima by abandoning older generations. \\ \midrule

Exact Learning to Branch (GNN-B\&B) & Iteratively predicts branching variable probabilities at each node of the B\&B tree. & State-of-the-art Predict-Many framework; replaces expensive branching calculations with GNN inference. \\ \midrule

Fast Iterative Localized Optimization (FILO) & Local search constrained by predefined spatial granularity and proximity limits. & Genuinely effective spatial heuristic designed for massive-scale routing problems. \\ \midrule

Firefly Algorithm (FA) & Models bioluminescent attraction where brightness exponentially decays over distance. & \textbf{Deceptive.} A minor, continuous local-search derivative of PSO; highly inefficient in complex combinatorial routing. \\ \midrule

Genetic Algorithm (GA) & Evolves candidate solutions using proportional selection, recombination (crossover), and mutation. & Foundational mathematical archetype for population-based search and schema theory. \\ \midrule

Genetic Programming Hyper-Heuristic (GPHH) & Evolves new heuristic rules or constructive policies using GP trees rather than searching solution space. & Powerful automated heuristic design paradigm; shifts the search space from solutions to algorithms. \\ \midrule

Great Deluge (GD) & Accepts moves better than a monotonically updating "water level" threshold. & Foundational deterministic alternative to SA; manages the "uphill" budget without stochasticity. \\ \midrule

Guided Indicators Hyper-Heuristic (GIHH) & Selects low-level routing heuristics dynamically, guided by performance indicators like convergence. & Effective multi-objective control mechanism; balances exploration without relying on thematic metaphor. \\ \midrule

Guided Local Search (GLS) & Escapes local optima by dynamically penalizing features via an augmented objective. & Foundational memory-based search trajectory adjustment mechanism. \\ \midrule

Harmony Search (HS) & Mimics musicians improvising to find acoustic harmony via pitch adjustments. & \textbf{Deceptive.} Mathematically proven to be a relabeled, constrained $(\mu + 1)$-Evolution Strategy. \\ \midrule

Hopfield Neural Network (HNN) & Iteratively refines a continuous state relaxation by minimizing an unsupervised energy landscape. & Foundational archetype of unsupervised neural optimization based on statistical physics. \\ \midrule

Hybrid Genetic Search (HGS) & Genetic Algorithm that applies local search to offspring prior to fitness evaluation. & State-of-the-art framework for routing optimization; highly transparent and rigorous engineering. \\ \midrule

Hybrid Genetic Search with Ruin and Recreate (HGS-RR) & HGS wrapper that utilizes advanced ruin and recreate principles instead of simple 2-opt. & Genuinely novel intersection of the premier population paradigm with state-of-the-art trajectory operators. \\ \midrule

Hybrid Volleyball Premier League (HVPL) & Merges the VPL competition framework with ACO routing and ALNS neighborhood search. & \textbf{Exposes the Metaphor:} VPL was too weak to solve routing; authors injected rigorous ACO/ALNS to achieve performance. \\ \midrule

Hyper-Heuristic Ant Colony Optimization (HH-ACO) & Uses a global pheromone matrix to probabilistically select which low-level heuristic operators to apply. & Legitimate abstraction; applies canonical ACO logic to navigate the heuristic operator space. \\ \midrule

Hyper-heuristic with Us and LK (HULK) & ACO-based selector managing Unstringing/Stringing and Lin-Kernighan operators. & Advanced adaptive operator management; avoids metaphor by explicitly stating hybrid components. \\ \midrule

Improving and Equal (IE) & Elitist strategy accepting candidates better than or equal to current solutions. & Codified gatekeeper for navigating high levels of neutrality in heuristic search spaces. \\ \midrule

Iterated Local Search (ILS) & Iterates local search by systematically applying perturbation operators to escape basins. & Foundational trajectory archetype utilizing intelligent mathematical restarts rather than randomization. \\ \midrule

ILS-RVND-SP & Uses ILS and randomized VNS to populate a pool of routes solved exactly via Set Partitioning. & Gold-standard matheuristic; explicitly combines rigorous local search with exact ILP models. \\ \midrule

K-Sparse Ant Colony Optimization (KS-ACO) & Modifies ACO pheromone updates to maintain a strictly $K$-sparse probability distribution. & Effective structural constraint preventing premature convergence; a rigorous mathematical modifier. \\ \midrule

Kernel Search (KS) & Decomposes the problem by intelligently restricting the active binary variables available to an iterative MILP solver. & Guarantees structural feasibility implicitly via the exact solver; superior to pure metaheuristics in tight polyhedra. \\ \midrule

Knowledge Guided Local Search (KGLS) & Uses data mining to group variables prior to deterministic local search. & Data-driven execution of GLS; precise, honest, and completely devoid of metaphor. \\ \midrule

Late Acceptance Hill-Climbing (LAHC) & Accepts candidates if they are better than or equal to a solution from $L$ iterations ago. & Superior scale-independent memory mechanism; allows plateau traversal without parameter recalibration. \\ \midrule

League Championship Algorithm (LCA) & Models a sports schedule where teams adapt "formations" based on win/loss fitness. & \textbf{Deceptive.} Functionally an exhaustive, steady-state Genetic Algorithm obscured by sports terminology. \\ \midrule

Learning-Augmented Column Generation & Uses Graph Neural Networks to predict strong branching outcomes, filtering candidates before MILP execution. & Traverses the B\&P search tree exponentially faster by replacing computationally exorbitant LP solver tests with ML inference. \\ \midrule

Lin-Kernighan-Helsgaun (LKH) & Searches for improving $k$-opt moves restricted by minimum spanning trees. & The pinnacle of TSP trajectory heuristics; purely mathematical and completely devoid of metaphor. \\ \midrule

Local Branching (VNSB) & Restricts the MILP search space by enforcing a strict mathematical bound on the Hamming distance from the incumbent. & Turns the exact solver into an aggressive neighborhood explorer, pruning millions of variables and nodes instantaneously. \\ \midrule

Local Search (LS) & Iteratively moves to a strictly improving neighboring solution until local optimum. & The absolute foundational trajectory-based optimization mechanism. \\ \midrule

Memetic Algorithm (MA) & Augments population-based search with intensive local trajectory refinement for individuals. & Foundational hybrid paradigm; mathematically bridges exploration (evolution) and exploitation (local optimization). \\ \midrule

Metropolis-Hastings (MH) & Accepts moves based on probability density ratios to sample from a distribution. & The core statistical mechanics foundation from physics that directly birthed Simulated Annealing. \\ \midrule

NeuRewriter & Uses an actor-critic RL framework to iteratively refine a solution via rewrite rules. & Directly maps classical trajectory search into an RL framework. \\ \midrule

Old Bachelor Acceptance (OBA) & Uses a non-monotone threshold that tightens on accept and loosens on reject. & Valid dynamic tuning schedule for SA/TA cooling pipelines; prevents search stagnation. \\ \midrule

Only Improving (OI) & The strictest form of elitist move acceptance, accepting moves only if $f(s') < f(s)$. & Purely exploitative gatekeeper; the core logic of standard greedy Hill Climbing. \\ \midrule

POPMUSIC & Dynamically partitions massive routing instances into subproblems optimized exactly by a BCP solver. & State-of-the-art decomposition framework; operates over massive neighborhoods to escape local optima efficiently. \\ \midrule

Particle Swarm Optimization (PSO) & Particles move guided by personal bests and global swarm knowledge. & Foundational continuous swarm archetype; heavily plagiarized by animal-based metaphors. \\ \midrule

Particle Swarm Optimization Memetic Algorithm (PSOMA) & Utilizes PSO as global wrapper, refining particle positions directly with trajectory local search. & Transparent combination of canonical PSO exploration with local search exploitation. \\ \midrule

Physics-Inspired GNN (PI-GNN) & Learns node assignments by directly minimizing physics-based Hamiltonian energy. & Modern Unsupervised Learning; replaces heuristic search with continuous energy gradients. \\ \midrule

Pointer Network (Ptr-Net) & Uses an LSTM with attention to output a permutation of graph nodes step-by-step. & Foundational SL model that birthed NCO; limited by reliance on optimal labels. \\ \midrule

Quantum Differential Evolution (QDE) & Replaces uniform random parameter distributions with chaotic probability functions. & Primarily a parameter tuning strategy; use of "Quantum" terminology is often more semantic than functional. \\ \midrule

Reactive Tabu Search (RTS) & Dynamically adjusts tabu list cardinality based on search history and cycle detection. & Foundational advancement in memory-based adaptive trajectory search; mathematically rigorous. \\ \midrule

Record-to-Record Travel (RRT) & Accepts candidates if they fall within a fixed scalar margin of the best-known record. & Deterministic acceptance rule; particularly effective for Cumulative VRPs where sequences are highly sensitive. \\ \midrule

Reinforcement Learning B\&B (RL-B\&B) & Interacts with the B\&B tree as an environment to minimize tree size via RL actions. & Rigorous integration of RL into exact mathematical programming spaces. \\ \midrule

Reinforcement Learning - Great Deluge HH (RL-GD-HH) & Uses RL to select low-level heuristics while GD acceptance governs state transitions. & Rigorous cross-domain hybridization; controls a thermodynamic criterion with learned policies. \\ \midrule

Relaxation Enforced Neighborhood Search (RENS) & Constructs a sub-MIP that represents the space of all feasible roundings of the LP optimum. & A powerful start heuristic that doesn't require an incumbent solution; mathematically proves if any rounding exists for a given LP point. \\ \midrule

Relaxation Induced Neighborhood Search (RINS) & Fixes binary variables matching the continuous LP relaxation and solves a restricted MILP on the unfixed variables. & Exploits commercial solvers as black-box heuristic operators; highly effective for rapidly improving integer bounds. \\ \midrule

Restricted Master Heuristics & Replaces exact ESPPRC pricing subproblems with fast metaheuristics to generate columns for a set partitioning MILP. & Overcomes the exponential column generation bottleneck in B\&P, rapidly establishing near-optimal continuous bounds. \\ \midrule

Sequence-based Selection HH (SS-HH) & Probabilistic SS-HH governs threshold transitions for Great Deluge local search. & Legitimate parameter control; applies reinforcement learning to guide a thermodynamic heuristic. \\ \midrule

Simulated Annealing (SA) & Probabilistically accepts worse moves based on Boltzmann distribution and temperature. & Foundational stochastic criterion; manages global exploration-exploitation but is sensitive to objective scale. \\ \midrule

Simulated Annealing Neighborhood Search (SANS) & Explores diverse neighborhood structures utilizing SA's cooling schedule. & Transparent hybridization; correctly attributes its mathematical foundations to SA/LS. \\ \midrule

Sine Cosine Algorithm (SCA) & Fluctuates agents toward/away from best solution using sine/cosine functions. & \textbf{Deceptive.} Meaningless trigonometric "math-washing" applied to standard vector translation operations. \\ \midrule

Slack Induction by String Removal (SISR) & Removes spatial sequences of variables and reinserts them to minimize slack. & Genuinely novel spatial meta-heuristic; achieved state-of-the-art results in large-scale routing. \\ \midrule

Smart Waste Collection - Two-Commodity Flow (SWC-TCF) & Uses two continuous flow variables to simultaneously represent waste load and empty capacity, dynamically determining the optimal bins to visit to maximize profit based on real-time sensor data. & Exact solver implementation; rigorously extends the CVRP two-commodity flow archetype to profitable routing under strict service level constraints. \\ \midrule

Soccer League Competition (SLC) & Simulates matches, substitutions, and standings to optimize real-valued networks. & \textbf{Deceptive.} Convoluted multi-population Genetic Algorithm burdened with unnecessary sports rule constraints. \\ \midrule

Step Counting Hill Climbing (SCHC) & Memory-based criterion where a cost bound remains fixed for a specific number of steps. & Scale-independent alternative to SA; offers high performance on examination timetabling benchmarks. \\ \midrule

Threshold Accepting (TA) & Accepts deteriorating moves if they fall within a predefined threshold range. & Deterministic alternative to SA; allows reorganization of sequences to find globally efficient tours. \\ \midrule

Variable Neighborhood Search (VNS) & Systematically switches neighborhood structures when trapped in local optima. & Foundational framework establishing rigorous rules for trajectory diversification. \\ \midrule

Vehicle Routing Problem with Profits Optimizer (VRPP) & Data-driven exact optimizer maximizing collection profits via a rolling horizon MILP. & Exact solver implementation; utilizes HMMs to handle uncertainty while proving global daily optimality. \\ \midrule

Volleyball Premier League (VPL) & Mimics volleyball coaching via substitution, coaching, and learning phases. & \textbf{Deceptive.} Fundamentally flawed GA/ES hybrid that relies on thematic phases rather than rigorous exploitation. \\ \bottomrule
\end{longtable}
\end{center}

\section{Appendix B: Algorithms Hyperparameters}
\vspace*{-25pt} % Space after section title
\begin{table}[H]
    \centering
    \begin{tabular}{|c||c|c|c|c|c|c|c|c|c|c|c|c|c|}
    \hline
    \multicolumn{14}{|c|}{Reinforcement Learning Hyperparameters}\\
    \hline
    Model & Seed & ND & BS & LR & HD & ED & $\lambda_{w}$ & $\lambda_{l}$ & $\lambda_{o}$ & $\tau$ & $\lambda_o$ & $\sigma$ & $\epsilon$\\ \hline
    AM & $42$  & $31$ & $256$ & $10^{-4}$ & $512$ & 128 & $0.5$ & $0.5$ & $0.5$ & - & - & $0.5$ & -\\
    AMGC & $42$ & $31$ & $256$ & $10^{-4}$ & $512$ & 128 & $0.5$ & $0.5$ & $0.5$ & - & - & - & -\\\hline
    \end{tabular}
    \vspace*{10pt} % Space between table and caption
    \caption{Hyperparameters used to train the various models during the reinforcement learning setting. We trained all models using the RMSProp optimizer with the default PyTorch parameters. NE is the number of training days (equivalent to epochs or the length of an episode), BS is the batch size, LR is the learning rate, HD is the hidden layers dimension, ED is the embedding layers dimension, $\lambda_{w}$ is the waste collected cost weight, $\lambda_{l}$ is the length traveled cost weight, $\lambda_{o}$ is the overflows cost weight, $\tau$ is the Softmax temperature, $\lambda_o$ is the odd-one-out network loss scale, $\sigma$ is the noise standard deviation and $\epsilon$ is the FGSM adversarial attack parameter.}
    \label{tab:hyperparams_reinforcement}
\end{table}
\end{document}
